\chapter{Implementation}
\label{ch:implementation}

\section{Language and build system}

The implementation is in C++20, compiled with GCC~14.2 on Linux (WSL2).
CMake~$\ge 3.20$ orchestrates the build. We deliberately keep external
dependencies minimal: GoogleTest is fetched at configure time via
\texttt{FetchContent}, the Python plotting tooling needs only
\texttt{pandas}, \texttt{matplotlib}, and \texttt{numpy}, and nothing
else is required on the path of the C++ binaries themselves. The choice
of C++ over Python was driven by the high asymptotic complexity of the
CWZ algorithm: even an $O(|V|^{4}|E|^{3})$ inner loop must run on graphs
with $|V|\approx 30$--$50$ to give meaningful empirical data, and
Python's interpretation overhead would limit the reachable size
considerably.

Edge weights are stored as 64-bit signed integers (\texttt{Weight =
std::int64\_t}). The algorithm's edge classification compares
\textit{exact} equalities like $d_S[u]+w(u,v)=d_S[v]$, which
floating-point arithmetic would render unreliable. Generators emit
weights uniformly drawn from $[1, 10]$.

\section{Module layout}

The source tree is split into small modules under \texttt{src/}:
\texttt{graph/} (a directed graph with adjacency lists indexed by stable
edge IDs and both out- and in-neighbour lists), \texttt{shortest\_path/}
(Dijkstra forward and reverse), \texttt{two\_vdp/} (2-VDP-in-DAG and a
brute-force oracle for it), \texttt{cwz/} (reductions and the layered
NSP algorithm), \texttt{baseline/} (brute force and Yen), and
\texttt{cli/} (a thin \texttt{nsp-cli} wrapper). Graph generators live
under \texttt{generators/}; unit tests under \texttt{tests/unit/};
correctness fuzzers under \texttt{tests/correctness/}; benchmarks and
plotting under \texttt{bench/}.

\section{Key data structures}

The \texttt{Graph} class assigns a stable identifier to each edge at
insertion time and exposes both \texttt{out\_edges(v)} and
\texttt{in\_edges(v)} lists of edge IDs. Maintaining the reverse
adjacency in parallel makes reverse Dijkstra a single pass over in-edges
rather than a transpose-then-search operation, and it lets the layered
algorithm classify each edge as forward/back-edge once and store the
result on the side.

For the residual subgraph $G[V\setminus U]$ used inside the layered
algorithm we never materialize the subgraph; instead we run a
vertex-masked Dijkstra that skips any edge whose endpoint lies in
$U$. This avoids allocating an $O(|V|+|E|)$ subgraph for each of the
many $(A,B)$ candidates explored.

\section{Paper-faithful pipeline}
\label{sec:pipeline}

The top-level \texttt{cwz\_nsp(g, s, t)} runs the paper's three-stage
pipeline:
\begin{enumerate}
  \item \texttt{reduce\_to\_straight(g, s, t)} (the paper's
    \textsc{NextSP}) folds every off-shortest-path vertex, replacing it
    with shortcut edges between its in- and out-neighbours, leaving an
    $(s,t)$-straight graph;
  \item \texttt{reduce\_to\_layered} (the paper's
    \textsc{NextSP-Straight}) turns the straight graph into a
    \emph{strictly} $(s,t)$-layered one and, as a by-product, records a
    set of candidate NSP paths (see below);
  \item \texttt{cwz\_nsp\_layered} (the paper's \textsc{NextSP-Layered})
    runs the Lemma~5.3 $6$-tuple enumeration on the strictly layered
    graph;
  \item the answer is the cheapest among the recorded candidates and the
    layered result, mapped back to original-graph edges via a
    \texttt{path\_expand[]} table that each reduction maintains.
\end{enumerate}

\paragraph{\textsc{NextSP-Straight} removes violating back-edges and
records candidates.} The subtle and important part is step~2. The
paper's $(s,t)$-layered definition (Definition~4.1) requires that every
back-edge go \emph{strictly backward} in layer. An $(s,t)$-straight
graph can contain back-edges that go forward (or sideways) in distance
-- an edge $(u,v)$ with $d_G(s,u)\le d_G(s,v)$ but
$d_G(s,u)+w(u,v)>d_G(s,v)$. \textsc{NextSP-Straight} \emph{removes} each
such edge, but first records the candidate path
$P_{s\to u}\circ(u,v)\circ P_{v\to t}$ (shortest forward paths to $u$
and from $v$), which the paper proves is always a valid simple
not-shortest path. Forward edges that skip layers are subdivided through
intermediate-distance vertices; strictly-backward back-edges are kept
unchanged. The result is a graph in which \emph{every} back-edge is
strictly backward -- exactly the structure Lemma~5.1 guarantees for an
$(s,t)$-layered graph, and exactly the structure the layered enumeration
relies on. Our implementation (\texttt{reduce\_to\_layered}) records the
removed-edge candidates in \texttt{ReductionResult::candidates}, and
\texttt{cwz\_nsp} folds them into the final minimum.

\section{The layered NSP algorithm}
\label{sec:layered}

\texttt{cwz\_nsp\_layered} is a direct transcription of
\textsc{NextSP-Layered} (paper Section~5). It enumerates every
$6$-tuple $(A,B,X',X,Y',Y)$ with $A,B$ incident to back-edges,
$\lambda(X')=\lambda(Y')=\lambda(X)-1=\lambda(Y)-1$, and (implied by the
layer constraints) $d(A)>d(B)$. For each tuple it finds the pair of
vertex-disjoint forward paths $P_1=s\to X'\to X\to A$ and
$P_2=B\to Y'\to Y\to t$ via two calls to the 2-VDP-in-DAG subroutine --
one for the layers below the barrier ($s\to X'$, $B\to Y'$) and one for
the layers above ($X\to A$, $Y\to t$), which are automatically disjoint
because forward paths are monotone in layer. It then completes the
candidate with a shortest $A\to B$ path in the residual graph
$G[(V\setminus(V(P_1)\cup V(P_2)))\cup\{A,B\}]$ and keeps the cheapest
result. The Key Lemma (Lemma~5.3) guarantees this finds an NSP whenever
one exists.

The boundary cases $A=X$ (the upper part of $P_1$ is empty) and $B=Y'$
(the lower part of $P_2$ is empty) are common on small graphs and must
be admitted: the layer constraint on $B$ is $\lambda(B)\le L$, not
$\lambda(B)<L$, and infeasible tuples are rejected by the 2-VDP calls
rather than excluded up front. Tightening either of these silently loses
NSPs whose first back-edge sits exactly one layer above its last.

\section{The 2-VDP subroutine}
\label{sec:twovdp}

Given a DAG $D=(V,E)$ and vertices $s_1,t_1,s_2,t_2$, the subroutine
must find paths $P_1\colon s_1\to t_1$ and $P_2\colon s_2\to t_2$ with
$V(P_1)\cap V(P_2)=\emptyset$, or report that no such pair exists. The
layered algorithm calls it on the subgraph of forward edges. That
subgraph is acyclic: a forward edge $(u,v)$ satisfies
$d(s,v)=d(s,u)+w(u,v)>d(s,u)$, since all weights are positive.

We implement the algorithm of Fortune, Hopcroft and
Wyllie~\cite{FortuneHopcroftWyllie1980} for two paths
(\texttt{two\_vdp\_in\_dag}). It runs in $O(|V|^{2}+|V||E|)$ time, which
is slower than Tholey's $O(|V|+|E|)$ algorithm~\cite{Tholey2012} but
much shorter to implement.

\paragraph{Algorithm.} If $s_1=s_2$, $t_1=t_2$, $s_1=t_2$ or $s_2=t_1$,
some vertex would have to lie on both paths, and the answer is
negative. Otherwise fix a topological order of $D$ and let $\pi(x)$ be
the position of vertex $x$. A \emph{state} is a pair $(u,v)$ of distinct
vertices, meaning that $P_1$ has been built from $s_1$ up to $u$ and
$P_2$ from $s_2$ up to $v$. In a state $(u,v)\neq(t_1,t_2)$ exactly one
of the two paths is extended:
\begin{itemize}
  \item $P_1$ is extended if $u\neq t_1$ and either $v=t_2$ or
    $\pi(u)<\pi(v)$; this gives the moves $(u,v)\to(w,v)$ for every edge
    $(u,w)\in E$ with $w\neq v$;
  \item otherwise $P_2$ is extended, giving the moves $(u,v)\to(u,w)$
    for every edge $(v,w)\in E$ with $w\neq u$.
\end{itemize}
In words, the path whose last vertex comes earlier in the topological
order is extended, and a path that has reached its target is never
extended again. If $P_2$ is extended then $v\neq t_2$: either $u\neq
t_1$ and the rule forces it, or $u=t_1$ and then $v\neq t_2$ because the
state is not $(t_1,t_2)$.

The moves form a directed graph $H$ on at most $|V|(|V|-1)$ states. The
algorithm runs a depth-first search in $H$ from $(s_1,s_2)$ and stores,
for every state, the edge of $D$ through which the state was first
reached. If $(t_1,t_2)$ is reached, the paths are recovered by following
the stored edges back to $(s_1,s_2)$. Each stored edge ends in exactly
one of the two current vertices, since $u\neq v$, and that tells us
whether it belongs to $P_1$ or to $P_2$. If the search ends without
reaching $(t_1,t_2)$, the answer is negative.

\paragraph{Correctness.} A walk in $H$ starting at $(s_1,s_2)$ records
two sequences of vertices of $D$, one for each path, by listing the
vertices added by the moves.

\begin{lemma}\label{lem:vdp-sound}
For every walk in $H$ from $(s_1,s_2)$ to a state $(u,v)$, the recorded
sequences are vertex-disjoint paths $P_1\colon s_1\to u$ and
$P_2\colon s_2\to v$.
\end{lemma}

\begin{proof}
Induction on the length of the walk. Before any move is made, the state
is $(s_1,s_2)$, $P_1$ consists of the single vertex $s_1$ and $P_2$ of
the single vertex $s_2$. They are disjoint because $s_1\neq s_2$. Suppose the claim holds in $(u,v)$ and
the next move extends $P_1$ along an edge $(u,w)$. Since $D$ is acyclic,
$\pi(w)>\pi(u)\ge\pi(x)$ for every $x\in V(P_1)$, so $P_1$ remains a
path. It remains to show $w\notin V(P_2)$. The move requires $w\neq v$.
Let $x$ be a vertex of $P_2$ other than $v$. Then $P_2$ has been
extended at least once; consider its last extension, a move
$(u',v')\to(u',v)$. The rule chose $P_2$ in state $(u',v')$, so either
$u'=t_1$, or $v'\neq t_2$ and $\pi(v')<\pi(u')$. If $u'=t_1$, then $P_1$
was never extended afterwards and $u=t_1$, which contradicts the fact
that $P_1$ is extended now. Hence $\pi(v')<\pi(u')$. The vertex $x$ lies
on $P_2$ no later than $v'$ and $u'$ lies on $P_1$ no later than $u$, so
\[
  \pi(x)\le\pi(v')<\pi(u')\le\pi(u)<\pi(w),
\]
and $x\neq w$. The case in which $P_2$ is extended is symmetric: the
last extension of $P_1$ took place in a state $(u',v')$ with $u'\neq
t_1$ and either $v'=t_2$ or $\pi(u')<\pi(v')$; the first option would
mean $v=t_2$, but $P_2$ is being extended, so
$\pi(x)\le\pi(u')<\pi(v')\le\pi(v)<\pi(w)$ for every $x\in V(P_1)$
other than $u$.
\end{proof}

\begin{lemma}\label{lem:vdp-complete}
If there are vertex-disjoint paths $Q_1\colon s_1\to t_1$ and
$Q_2\colon s_2\to t_2$ in $D$, then $(t_1,t_2)$ is reachable from
$(s_1,s_2)$ in $H$.
\end{lemma}

\begin{proof}
We build a walk in $H$ that keeps $u$ on $Q_1$ and $v$ on $Q_2$. This
holds in $(s_1,s_2)$. In a state $(u,v)\neq(t_1,t_2)$ the rule selects a
path that has not reached its target, as observed above. Move its last
vertex to the next vertex $w$ of $Q_1$ or $Q_2$ respectively. Since
$Q_1$ and $Q_2$ are disjoint, $w$ differs from the last vertex of the
other path, so this is a move of $H$, and the invariant still holds.
Every move advances along $Q_1$ or $Q_2$, so after
$|E(Q_1)|+|E(Q_2)|$ moves the walk is in $(t_1,t_2)$.
\end{proof}

The depth-first search visits every state reachable from $(s_1,s_2)$,
so by Lemma~\ref{lem:vdp-complete} it finds $(t_1,t_2)$ whenever a
solution exists. The stored edges describe one walk from $(s_1,s_2)$ to
$(t_1,t_2)$, and by Lemma~\ref{lem:vdp-sound} the paths recovered from
it are vertex-disjoint. Hence \texttt{two\_vdp\_in\_dag} returns a pair
of paths if and only if one exists, and every returned pair is correct.

\paragraph{Running time.} The topological order is computed with
Kahn's algorithm in $O(|V|+|E|)$ time; if the input contains a cycle,
the order is incomplete and the function throws instead of answering.
Each state is expanded at most once. Expanding $(u,v)$ scans the
out-edges of one of $u$ and $v$, so the total work of the search is at
most
\[
  \sum_{u\in V}\sum_{v\in V}\bigl(\deg^{+}(u)+\deg^{+}(v)\bigr)
  = 2|V||E|.
\]
Together with the $|V|^{2}$-entry array of stored edges this gives
$O(|V|^{2}+|V||E|)$ time and $O(|V|^{2})$ memory per call.

\section{Testing strategy}
\label{sec:correctness}

The project ships:
\begin{itemize}
  \item $110$ unit tests covering the graph data structure, both
    Dijkstras, 2-VDP ($54$ tests), both
    reductions (including that \textsc{NextSP-Straight} removes
    forward/sideways back-edges and records the right candidates), the
    layered NSP algorithm on hand-built examples, both baselines, and
    the end-to-end \texttt{cwz\_nsp} pipeline. All pass.
  \item A brute-force 2-VDP solver
    (\texttt{src/two\_vdp/brute\_force\_two\_vdp}) used as an oracle for
    \texttt{two\_vdp\_in\_dag}. It enumerates all simple $s_1\to t_1$
    and $s_2\to t_2$ paths by depth-first search, stores the vertex set
    of each path as a $64$-bit mask, and reports a solution if and only
    if some $s_1\to t_1$ mask and some $s_2\to t_2$ mask have an empty
    intersection. It takes exponential time and is limited to $64$
    vertices, but it shares no code with the Fortune--Hopcroft--Wyllie
    implementation and does not rely on acyclicity. The unit tests
    compare the two on every terminal quadruple $(s_1,t_1,s_2,t_2)$ of
    all DAGs on $4$ labelled vertices ($393{,}216$ queries), of all
    $1{,}024$ edge subsets of a $5$-vertex transitive tournament under a
    random relabelling ($640{,}000$ queries), and of $400$ random DAGs
    with up to $7$ vertices and occasional parallel edges ($266{,}533$
    queries). Every pair of paths returned by \texttt{two\_vdp\_in\_dag}
    is also checked to consist of two vertex-disjoint paths with the
    right endpoints. The two solvers agree on all $1{,}299{,}749$
    queries. Hand-built tests cover the degenerate cases $s_i=t_i$,
    parallel edges, rejection of cyclic inputs, and instances in which
    only the crossed pairing $s_1\to t_2$, $s_2\to t_1$ is routable.
  \item A property-based fuzz \texttt{fuzz\_brute\_vs\_yen} that
    confirms brute force and Yen's NSP heuristic agree on cost across
    thousands of random small graphs (mismatches would indicate a bug
    in either side).
  \item A differential fuzz \texttt{fuzz\_cwz\_vs\_brute} that confirms
    the full CWZ pipeline matches brute force. We measured \emph{zero}
    mismatches across $34{,}000$ random trials drawn from $35$ distinct
    RNG seeds: $|V|\le 10$ over $20$ seeds ($20{,}000$ trials),
    $|V|\le 12$ over $10$ seeds ($10{,}000$ trials), and $|V|\le 14$
    over $5$ seeds ($4{,}000$ trials). Of these, $23{,}624$ instances
    had an NSP and the remaining $10{,}376$ had none; both algorithms
    agreed on which case held every time. The seeds are consecutive
    integers ($1001$--$1020$, $2001$--$2010$, $3001$--$3005$), so the
    campaign is reproducible verbatim.
\end{itemize}

The combined test suite establishes correctness on small graphs;
agreement with brute force on all $1200$ benchmark instances
(Section~\ref{sec:agreement}) extends the evidence to larger graphs.

\section{The CLI}

\texttt{nsp-cli} reads a graph from a DIMACS-style text file (first line
$n\;m$, then $m$ lines of $u\;v\;w$) and runs the selected algorithm.
Source and sink default to vertex $0$ and vertex $n-1$. The output is
one \texttt{nsp\_cost} line and a list of edges describing the recovered
NSP.
